\documentclass[twoside,10pt]{article}
\usepackage[T1]{fontenc}
\usepackage{lmodern}
\usepackage{amsmath,amssymb,amsthm,amsfonts,mathtools}
\usepackage[letterpaper,total={5.5in,8in},vmarginratio=1:1,hmarginratio=1:1,headsep=21pt]{geometry}
\usepackage{xcolor}
\usepackage[breaklinks=true,linktocpage=true]{hyperref}
\hypersetup{colorlinks=true,linkcolor=purple,urlcolor=purple,citecolor=blue!80!black,
pdfauthor={Leslie P. Polzer},pdftitle={Coefficient criteria for finite nonnegative linearization},
pdfsubject={Recurrence conditions, quadratic systems and quantitative stability for finite orthogonal polynomials},
pdfkeywords={orthogonal polynomials, Krawtchouk polynomials, nonnegative linearization, quadratic transformation, hypergroups, Lean}}
\medmuskip=4mu
\thickmuskip=8mu
\setlength{\emergencystretch}{2em}
\numberwithin{equation}{section}
\newtheoremstyle{slbody}{8pt}{8pt}{\normalfont\slshape}{}{\bfseries}{.}{0.5em}
{\thmname{#1}\thmnumber{ #2}\thmnote{ {\normalfont\itshape(\ignorespaces#3\unskip)}}}
\theoremstyle{slbody}
\newtheorem{theorem}{Theorem}[section]
\newtheorem{lemma}[theorem]{Lemma}
\newtheorem{corollary}[theorem]{Corollary}
\newtheorem{proposition}[theorem]{Proposition}
\theoremstyle{definition}
\newtheorem{definition}[theorem]{Definition}
\theoremstyle{remark}
\newtheorem{remark}[theorem]{Remark}
\renewcommand{\qedsymbol}{$\blacksquare$}
\newcommand{\qint}[1]{[#1]_q}
\title{Coefficient criteria for finite nonnegative linearization}
\author{Leslie P. Polzer}
\newcommand{\authoraffiliation}{Independent Researcher}
\newcommand{\authoremail}{polzer@fastmail.com}
\date{October 19, 2026}
\makeatletter
\renewcommand{\section}{\@startsection{section}{1}{\z@}{-3.5ex \@plus -1ex \@minus -.2ex}{2.3ex \@plus .2ex}{\normalfont\large\bfseries}}
\renewcommand{\@seccntformat}[1]{\csname the#1\endcsname.\enspace}
\def\tagform@#1{\maketag@@@{(\ignorespaces{\oldstylenums{#1}}\unskip\@@italiccorr)}}
\renewcommand{\eqref}[1]{\textup{{\normalfont(\oldstylenums{\ref{#1}}}\normalfont)}}
\newcommand{\shorttitle}{FINITE NONNEGATIVE LINEARIZATION}
\newcommand{\shortauthors}{LESLIE P. POLZER}
\newcommand{\ps@bookheader}{%
\renewcommand\@oddfoot{\hfil}\renewcommand\@evenfoot{\hfil}%
\renewcommand\@oddhead{\ifnum\value{page}>1{\small\hfil\shortauthors}\hfil\thepage\else\hfil\fi}%
\renewcommand\@evenhead{\thepage\hfil\small\shorttitle\hfil}}
\pagestyle{bookheader}
\renewcommand{\maketitle}{\begin{center}{\large\bfseries\@title}\\\vskip36pt
{\scshape\@author}\par\medskip{\small\authoraffiliation}\par\smallskip
{\small\texttt{\authoremail}}\par\smallskip{\small\@date}\par\vskip36pt\end{center}}
\makeatother
\renewenvironment{abstract}{\quotation\noindent\small{\bfseries\abstractname.\enspace}}{\endquotation}
\begin{document}
\maketitle

\begin{abstract}
We give two sufficient criteria stated directly in the coefficients of a finite stochastic three-term recurrence. The first assumes that the sum of its diagonal and downward coefficients has positive geometric increments with ratio $q\geq1$, and that the ratio of those coefficients is constant and nonnegative. A finite algebra construction gives nonnegative linearization, including products beyond the terminating degree. This includes the full positive Krawtchouk range and the dual $q$-Krawtchouk family with nonconstant increments. The underlying quadratic algebra occurs in Koelink's work on type B Hecke algebras. The second criterion gives explicit perturbation inequalities admitting a relatively open neighborhood of recurrence arrays around every Krawtchouk system in the interior of the positive linearization range. For Szwarc's symmetric quadratic system of length $2N+1$, a graded version of the algebra proves nonnegative linearization for $p_0\leq1/2$, and a third criterion admits explicit symmetric neighborhoods of it. None of the criteria requires computing general linearization coefficients.
\vskip5pt
\noindent\textbf{Keywords.}\enspace Finite orthogonal polynomials; nonnegative linearization; Krawtchouk polynomials; dual $q$-Krawtchouk polynomials; quadratic transformation; polynomial hypergroups.
\vskip5pt
\noindent\textbf{MSC2020 Classification.}\enspace 33C45; 33D45; 42C05.
\end{abstract}

\vskip50pt
\baselineskip=13.5pt
\section{Introduction}\label{sec:introduction}

Nonnegative linearization (property (P)) asks whether the product of two orthogonal polynomials is a nonnegative linear combination of polynomials in the same basis. For a finite system, the identity is understood on the orthogonality support. This distinction is essential when the formal product degree exceeds the terminal degree. In a normalization with $R_i(1)=1$, nonnegative coefficients have sum one and define the product probabilities of a finite polynomial hypergroup.

The Krawtchouk case is classical. Eagleson~\cite{eagleson} studied positive definite sequences for these polynomials in 1969; Szwarc~\cite[\S5, Problem~2]{szwarc} records the nonnegative-linearization range as $0<p_0\leq1/2$ for binomial weights $\binom Njp_0^{N-j}(1-p_0)^j$ when normalized at $j=0$.\footnote{The linked author version, \S5, Problem~2, printed p.~125, uses inconsistent parameter conventions: its displayed weights have mean $N(1-p)$, whereas its degree-zero recurrence gives $K_1(x)=(x-pN)/((1-p)N)$. Thus the displayed weights give $\int K_1\,d\mu=(1-2p)/(1-p)$. Corollary~\ref{cor:kraw} derives our weights and normalization directly.} Our proof below fixes that convention directly from a two-point model. It does not treat Krawtchouk positivity itself as a new result.

Recurrence methods also have a long history. In the second note of his 1969 report, Askey~\cite[Part~II, Theorem~1 and Corollary~1]{askey} gives finite recurrence conditions leading to a hyperbolic difference-equation maximum principle. The stated polynomial conclusion concerns products with degree sum at most $N$. Szwarc's later characterization~\cite[Theorem~1]{szwarc2003} treats finite as well as infinite systems through an auxiliary discrete hyperbolic problem. His 2005 survey nevertheless asks for criteria for finite systems, including Krawtchouk and its quadratic system $Q_n$ (Section~\ref{sec:szwarc})~\cite[\S5, Problem~2]{szwarc}. Here we seek structural assumptions on the recurrence coefficients that imply positivity for every finite product.

Koelink~\cite[Theorems~3.1(ii) and 7.8]{koelink} constructed a triangular quadratic algebra associated with the type B Hecke algebra, and identified the corresponding $q$-Krawtchouk spherical functions. The dual recurrence in that construction is the geometric family used below. Our algebraic proof extracts its positive multiplication mechanism and expresses the resulting sufficient condition as two scalar patterns in a stochastic recurrence: a constant ratio and geometric increments. No priority is claimed for the algebra or for the resulting named family. For comparison, Miclo~\cite[Proposition~3.3]{miclo} gives a spectral characterization of the hypergroup property at a boundary point for irreducible finite birth--death chains by entrywise nonnegativity of products of transition matrices shifted by eigenvalues of their leading principal submatrices.

The second criterion is quantitative stability. Around each Krawtchouk array in the interior of the positive linearization range ($0<r<1$, equivalently $0<p_0<1/2$), we exhibit an explicit radius in the space of recurrence coefficients within which all finite products remain nonnegative. Its proof uses one positive term to obtain a reference margin, an error recurrence to control perturbations, and orthogonality to preserve exact zeros outside the triangle support. The resulting neighborhoods contain independently varying coefficient arrays, rather than only a reparametrized classical family. The radius is deliberately conservative and depends on $N$; no uniform stability radius at the symmetric endpoint is asserted.

Section~\ref{sec:criterion} states the geometric criterion. Sections~\ref{sec:algebra} and~\ref{sec:normalization} give its algebraic proof. Section~\ref{sec:perturbation} proves the perturbation inequalities, including explicit nongeometric examples. Section~\ref{sec:applications} fixes the Krawtchouk normalization and compares the geometric recurrence with Koelink's construction. The final section lists questions left by these sufficient conditions.

\section{Statement and conventions}\label{sec:criterion}
For integers $j\geq0$ put $[j]_q=1+q+\cdots+q^{j-1}$, with $[0]_q=0$; in particular $[j]_1=j$.
Consider real coefficients indexed by $0\leq i\leq N$, where $N\geq1$, satisfying
\begin{equation}\label{eq:basic}
\begin{gathered}
a_0=1,\quad b_0=c_0=0,\quad a_N=0,\quad a_i+b_i+c_i=1,\\
a_i>0\ (i<N),\qquad c_i>0\ (i>0).
\end{gathered}
\end{equation}
Set $R_{-1}=0$, $R_0=1$, and determine $R_1,\ldots,R_N$ by
\begin{equation}\label{eq:recurrence}
xR_i=a_iR_{i+1}+b_iR_i+c_iR_{i-1}\qquad(0\leq i<N).
\end{equation}
At $i=0$, \eqref{eq:recurrence} gives $R_1=x$; induction at $x=1$, using the row sums and $a_i>0$ for $i<N$, gives $R_i(1)=1$ for $0\leq i\leq N$.
The terminal relation is $(x-b_N)R_N-c_NR_{N-1}=0$ on the orthogonality support. Linearization means equality of functions on this support, rather than an unrestricted polynomial identity when the degree sum exceeds $N$.

\begin{samepage}
\begin{theorem}[Geometric increment criterion]\label{thm:criterion}
Assume \eqref{eq:basic}. Set $d_i=b_i+c_i$. Suppose
\begin{enumerate}
\item $b_i/c_i=\beta$ for every $1\leq i\leq N$, where $\beta\geq0$;
\item $\delta_i=d_i-d_{i-1}>0$ for $1\leq i\leq N$;
\item for some $q\geq1$, $\delta_{i+1}=q\delta_i$ whenever $1\leq i<N$.
\end{enumerate}
Then the recurrence has a positive orthogonality measure on exactly $N+1$ distinct real points. Its normalized polynomials satisfy
\begin{equation}\label{eq:positive}
R_mR_n=\sum_{k=0}^N g(m,n;k)R_k,
\qquad g(m,n;k)\geq0,\qquad\sum_{k=0}^N g(m,n;k)=1
\end{equation}
on that support for all $0\leq m,n\leq N$.
\end{theorem}
\end{samepage}

\begin{lemma}[Coefficient parametrization]\label{lem:parameters}
The hypotheses of Theorem~\ref{thm:criterion} are equivalent to
\begin{equation}\label{eq:parameters}
a_i=\frac{q^i[N-i]_q}{[N]_q},\qquad
b_i=(1-r)\frac{[i]_q}{[N]_q},\qquad
c_i=r\frac{[i]_q}{[N]_q}
\end{equation}
for $q\geq1$, $0<r\leq1$. When $N=1$, one may choose $q=1$.
\end{lemma}
\begin{proof}
The terminal sum gives $d_N=1$, and $d_0=0$. Geometric increments give $d_i=\delta_1[i]_q$, hence $\delta_1=1/[N]_q$. The constant ratio gives $c_i=d_i/(1+\beta)$ and $b_i=\beta d_i/(1+\beta)$. Put $r=(1+\beta)^{-1}$. Finally $[N]_q-[i]_q=q^i[N-i]_q$. These calculations reverse to give the converse, including $q=1$ and $N=1$.
\end{proof}

The parameter conventions used below are
\[
\begin{array}{c|l}
q & q\geq1\quad\text{(geometric increment ratio)}\\
\beta & b_i/c_i\geq0\\
r & 1/(1+\beta)\in(0,1]\\
p & 1/r=1+\beta\geq1\\
p_0 & r/(1+r)\in(0,1/2]\\
\text{symmetric Krawtchouk case} & r=1,\quad p_0=1/2
\end{array}
\]

\section{A finite algebra with positive multiplication}\label{sec:algebra}
Fix $p\geq1$ and $q\geq1$. Starting with $A_{N+1}=\mathbb R$, form recursively the commutative real algebras
\begin{equation}\label{eq:algebra}
A_j=A_{j+1}[z_j]\big/
\left(z_j^2-\left((q-1)\sum_{l=j+1}^N z_l+p-1\right)z_j-pq^{N-j}\right).
\end{equation}
Write $A=A_1$. For $S\subseteq\{1,\ldots,N\}$ let $z_S=\prod_{j\in S}z_j$, with $z_\varnothing=1$.

\begin{lemma}[Basis and positivity]\label{lem:basis}
The elements $z_S$ form a basis of $A$, and all their multiplication coefficients are nonnegative.
\end{lemma}
\begin{proof}
Division by a monic quadratic over any commutative coefficient ring gives a unique remainder of degree at most one. Applying this successively in \eqref{eq:algebra} gives the asserted basis and dimension $2^N$.
For positivity, replace a square by
\[
z_j^2=(q-1)\sum_{l>j}z_jz_l+(p-1)z_j+pq^{N-j}.
\]
Every displayed coefficient is nonnegative. This rewriting terminates: give $z_j$ weight $2^{N-j}$ and give a monomial the sum of the weights of its factors, counted with multiplicity. Each replacement strictly decreases this nonnegative integer weight, because $l>j$ implies $2^{N-l}<2^{N-j}$. The final expression is in the unique squarefree basis, proving the claim.
\end{proof}

\begin{lemma}[Positive associative inner product]\label{lem:trace}
Let $\tau$ extract the coefficient of $z_\varnothing$. Then
\begin{equation}\label{eq:trace}
\tau(z_Sz_T)=\boldsymbol1_{S=T}\,p^{|S|}q^{\sum_{j\in S}(N-j)}.
\end{equation}
Consequently $\langle u,v\rangle=\tau(uv)$ is positive definite, and multiplication by every element of $A$ is self-adjoint.
\end{lemma}
\begin{proof}
Induct downwards through the tower. If $j$ belongs to exactly one of $S,T$, their product has the form $z_jv$ with $v\in A_{j+1}$ and has zero constant coefficient in $z_j$. If it belongs to both, the constant coefficient in $z_j$ is $pq^{N-j}z_{S\setminus\{j\}}z_{T\setminus\{j\}}$ by \eqref{eq:algebra}. If it belongs to neither, pass to $A_{j+1}$. Iteration proves \eqref{eq:trace}. All diagonal factors are positive. Associativity and commutativity give $\tau((hu)v)=\tau(u(hv))$.
\end{proof}

Put $E_k=\sum_{|S|=k}z_S$ for $0\leq k\leq N$, and set $E_{-1}=E_{N+1}=0$.

\begin{lemma}[Three-term multiplication]\label{lem:three}
For $0\leq k\leq N$,
\begin{equation}\label{eq:three}
E_1E_k=[k+1]_qE_{k+1}+(p-1)[k]_qE_k+
pq^{k-1}[N-k+1]_qE_{k-1}.
\end{equation}
The last term is interpreted as zero for $k=0$.
\end{lemma}
\begin{proof}
For $k=0$ the identity reads $E_1=E_1$. For $N=1,k=1$ it is exactly $z_1^2=(p-1)z_1+p$.
For the induction on $N$, write $F_j$ for the elementary squarefree sums in $z_2,\ldots,z_N$, with $F_j=0$ outside $0\leq j\leq N-1$. Then $E_k=F_k+z_1F_{k-1}$ and $E_1=F_1+z_1$. For $1\leq k\leq N-1$, use \eqref{eq:algebra} to obtain
\[
E_1E_k=F_1F_k+z_1F_k+qz_1F_1F_{k-1}
 +(p-1)z_1F_{k-1}+pq^{N-1}F_{k-1}.
\]
Apply the induction hypothesis to $F_1F_k$ and $F_1F_{k-1}$. The coefficients of $F_{k+1}$ and $z_1F_k$ are respectively $[k+1]_q$ and $1+q[k]_q=[k+1]_q$. The coefficients of $F_k$ and $z_1F_{k-1}$ are both $(p-1)[k]_q$, since $1+q[k-1]_q=[k]_q$. The coefficient of $F_{k-1}$ is
\[
pq^{k-1}[N-k]_q+pq^{N-1}=pq^{k-1}[N-k+1]_q,
\]
which is also the coefficient of $z_1F_{k-2}$. Regrouping gives \eqref{eq:three}.
For the endpoint $k=N$, $E_N=z_1F_{N-1}$, so the same square relation gives
\[
E_1E_N=qz_1F_1F_{N-1}+(p-1)z_1F_{N-1}+pq^{N-1}F_{N-1}.
\]
The tail endpoint identity is
\[
F_1F_{N-1}=(p-1)[N-1]_qF_{N-1}+pq^{N-2}F_{N-2}.
\]
Substitution gives $(p-1)[N]_qE_N+pq^{N-1}E_{N-1}$, as required. This separately checks the terminating boundary.
\end{proof}

\begin{proposition}[Positive subalgebra]\label{prop:subalgebra}
The space $B=\operatorname{span}\{E_0,\ldots,E_N\}$ is the subalgebra generated by $E_1$. Every product $E_mE_n$ has nonnegative coefficients in this basis.
\end{proposition}
\begin{proof}
All $E_k$ are independent, since their squarefree supports are nonempty and disjoint. Equation~\eqref{eq:three} shows that $B$ is invariant under multiplication by $E_1$. Since $[k+1]_q>0$, its interior recurrence also expresses each $E_{k+1}$ as a polynomial in $E_1$. Thus $B=\mathbb R[E_1]$, proving closure under multiplication, including products of total index greater than $N$.
Expand such a product as $\sum_k h_kE_k$. Its coefficient at each squarefree monomial $z_S$ is $h_{|S|}$. Lemma~\ref{lem:basis} makes that coefficient nonnegative. There is a subset of each cardinality $0,\ldots,N$, so every $h_k$ is nonnegative.
\end{proof}

\section{Normalization and the finite orthogonal system}\label{sec:normalization}
\begin{lemma}[Positive normalization]\label{lem:normalization}
Evaluation $z_j\mapsto pq^{N-j}$ defines an algebra character $\epsilon$ of $A$. Put $t_k=\epsilon(E_k)>0$. Then
\begin{equation}\label{eq:ratio}
t_0=1,\quad t_1=p[N]_q,\qquad
\frac{t_{k+1}}{t_k}=pq^k\frac{[N-k]_q}{[k+1]_q}\quad(0\leq k<N).
\end{equation}
\end{lemma}
\begin{proof}
For $h=N-j$, substitution in the right-hand side of \eqref{eq:algebra} gives
\[
(q-1)p^2q^h[h]_q+(p-1)pq^h+pq^h
=p^2q^h\big((q-1)[h]_q+1\big)=p^2q^{2h}.
\]
Thus all defining relations hold. The generating polynomial of the $t_k$ is
$T(u)=\prod_{j=0}^{N-1}(1+pq^ju)$, whose coefficients are strictly positive.
When $q>1$, the identity $(1+pu)T(qu)=(1+pq^Nu)T(u)$ gives
$(q^k-1)t_k=p(q^N-q^{k-1})t_{k-1}$ for $1\leq k\leq N$, which is \eqref{eq:ratio}. When $q=1$, $t_k=\binom Nk p^k$ and the same ratio follows directly.
\end{proof}

\begin{proof}[Proof of Theorem~\ref{thm:criterion}]
Use Lemma~\ref{lem:parameters} and put $p=1/r\geq1$. Construct $A,B$ as above and set $X=E_1/t_1$, $U_k=E_k/t_k$. Divide \eqref{eq:three} by $t_1t_k$ and use \eqref{eq:ratio}. The three coefficients are
\begin{align*}
\frac{[k+1]_qt_{k+1}}{p[N]_qt_k}
&=\frac{q^k[N-k]_q}{[N]_q},\\
\frac{(p-1)[k]_q}{p[N]_q}&=(1-r)\frac{[k]_q}{[N]_q},\\
\frac{pq^{k-1}[N-k+1]_qt_{k-1}}{p[N]_qt_k}
&=r\frac{[k]_q}{[N]_q}.
\end{align*}
At $k=0$ omit the third expression; at $k=N$ omit the first. Consequently $U_k=R_k(X)$ for every $k\leq N$, with the stated terminal relation.

The orthogonality support follows from the real spectral theorem in finite dimensions. On $B$ use the positive inner product from Lemma~\ref{lem:trace}. Multiplication by $X$ is self-adjoint. The vector $1=E_0$ is cyclic because the $E_k$ are polynomials of successive degrees in $X$ and form a basis. In an orthogonal eigenspace decomposition of a self-adjoint operator, the cyclic span of one vector has dimension at most the number of distinct eigenvalues. Here the cyclic span has dimension $N+1$, so there are exactly $N+1$ distinct eigenvalues, and the projection of $1$ on every eigenspace is nonzero. Give each eigenvalue the squared norm of that projection as its mass. All masses are positive, sum to $\langle1,1\rangle=1$, and yield
\[
\int f(x)h(x)\,d\mu(x)=\langle f(X),h(X)\rangle
\]
for real polynomials $f,h$, identifying an algebra element with its action on $1$. The $E_k$ are orthogonal by \eqref{eq:trace}, so the $R_k$ are orthogonal for this measure.

Proposition~\ref{prop:subalgebra} and positive $t_k$ give nonnegative product coefficients for the $U_k$. Applying the character $\epsilon$ gives their sum equal to one because $\epsilon(U_k)=1$. Polynomial identities in the algebra hold on every spectral point; conversely the $R_k$, having successive degrees at most $N$, form a basis of functions on the $N+1$ distinct points. This proves uniqueness and \eqref{eq:positive} for all products.
\end{proof}

\section{A criterion for small coefficient perturbations}\label{sec:perturbation}

Write the tridiagonal recurrence matrix with the row convention
\[
J_{i,i+1}=a_i,\qquad J_{i,i}=b_i,\qquad J_{i,i-1}=c_i,
\]
with all other entries zero and indices outside $0,\ldots,N$ omitted.

\begin{lemma}[Finite recurrence realization and triangle support]\label{lem:finite}
Every recurrence satisfying \eqref{eq:basic} has a positive orthogonality measure on exactly $N+1$ real points. If $J$ is its tridiagonal recurrence matrix, then its finite product coefficients satisfy
\begin{equation}\label{eq:matrix-product}
g(m,n;k)=[R_m(J)]_{n,k}.
\end{equation}
They vanish unless $|m-n|\leq k\leq\min(m+n,N)$, regardless of their signs inside this interval.
\end{lemma}
\begin{proof}
Choose positive $\pi_i$ with $\pi_0=1$ and $\pi_{i+1}=\pi_i a_i/c_{i+1}$. For $D=\operatorname{diag}(\sqrt{\pi_i})$, $H=DJD^{-1}$ is real symmetric. Induction in \eqref{eq:recurrence} gives the row-vector identity $e_0^\mathsf T R_i(J)=e_i^\mathsf T$, hence $e_0^\mathsf T R_i(H)=\pi_i^{-1/2}e_i^\mathsf T$. In particular $e_0$ is cyclic for $H$. The same finite spectral argument as in Theorem~\ref{thm:criterion} gives $N+1$ distinct eigenvalues and positive spectral masses at $e_0$, and
\[
\int R_iR_j\,d\mu=e_0^\mathsf TR_i(H)R_j(H)e_0
=\boldsymbol1_{i=j}\pi_i^{-1}.
\]
The terminal relation holds on this support: its polynomial $Q(x)=(x-b_N)R_N-c_NR_{N-1}$ satisfies $e_0^\mathsf TQ(J)=0$. Since $e_0$ is cyclic and $Q(J)$ commutes with $J$, $Q(J)=0$.
Multiplying a basis polynomial by $x$ acts by the corresponding row of $J$, including the terminal row. Induction on powers of $x$ therefore gives \eqref{eq:matrix-product}. Orthogonality also gives
\[
g(m,n;k)=\frac{\int R_mR_nR_k\,d\mu}{\int R_k^2\,d\mu}.
\]
If one of $m,n,k$ exceeds the sum of the other two, its polynomial is orthogonal to their product, whose degree is smaller. This proves the asserted support.
\end{proof}

\begin{lemma}[A uniform positive margin in the asymmetric case]\label{lem:margin}
For the Krawtchouk recurrence \eqref{eq:parameters} with $q=1$ and $0<r<1$, put $s=\min(r,1-r)$. Every coefficient in its triangle support obeys
\begin{equation}\label{eq:margin}
g(m,n;k)\geq (s/4)^N\quad\text{if }|m-n|\leq k\leq\min(m+n,N).
\end{equation}
\end{lemma}
\begin{proof}
Specialize the algebra above to $q=1$, and put $y_i=z_i/p$, where $p=1/r$. Then $y_i^2=(1-r)y_i+r$, and $U_m=\binom Nm^{-1}\sum_{|S|=m}y_S$. The reduction of a product $y_Sy_T$ has nonnegative coefficients, each chosen term being a product of at most $N$ factors from $\{r,1-r\}$. Its coefficient is therefore at least $s^N$.
For each admissible $k$, at least one such term has degree $k$. Assume $m\geq n$. If $k\leq m$, take $T\subseteq S$, with sizes $n,m$, and retain the linear term for exactly $k-(m-n)$ overlap coordinates. This number lies between $0$ and $n$. If $k\geq m$, choose sets with overlap $m+n-k$ and union of size $k\leq N$, and retain the linear term at every overlap coordinate.
Permutation symmetry makes the coefficient of $y_U$ depend only on $|U|=k$. Hence the coefficient in the normalized elementary-symmetric basis is at least
\[
\frac{s^N\binom Nk}{\binom Nm\binom Nn}\geq\frac{s^N}{4^N}.
\]
This uses a single positive term and a denominator bound, not an explicit formula for the product coefficients.
\end{proof}

\begin{theorem}[Explicit perturbation inequalities]\label{thm:perturbation}
Let $N\geq1$, $0<r<1$, $s=\min(r,1-r)$, and
\begin{equation}\label{eq:radius}
\eta_N(r)=\frac12\left(\frac{s}{32N}\right)^N.
\end{equation}
Suppose coefficients $\widetilde a_i,\widetilde b_i,\widetilde c_i$ satisfy \eqref{eq:basic}, with $\widetilde b_i\geq0$ for every $i$, and
\begin{equation}\label{eq:closeness}
\max_{0\leq i\leq N}\left\{
\left|\widetilde a_i-\frac{N-i}{N}\right|,
\left|\widetilde b_i-\frac{(1-r)i}{N}\right|,
\left|\widetilde c_i-\frac{ri}{N}\right|\right\}\leq\eta_N(r).
\end{equation}
Then the finite system has nonnegative linearization. Throughout its triangle support,
\[
\widetilde g(m,n;k)\geq\tfrac12(s/4)^N.
\]
The product coefficients sum to one.
\end{theorem}
\begin{proof}
Use tildes for the perturbed system and untilded letters for its $q=1$ reference at $r$. Let $\eta=\eta_N(r)$ and use the maximum absolute row-sum matrix norm. Both $J$ and $\widetilde J$ are stochastic matrices, so $\|J\|_\infty=\|\widetilde J\|_\infty=1$ and $\|\widetilde J-J\|_\infty\leq3\eta$. Put $B_m=R_m(J)$, $\widetilde B_m=\widetilde R_m(\widetilde J)$ and $e_m=\|\widetilde B_m-B_m\|_\infty$. Theorem~\ref{thm:criterion} at $q=1$ and \eqref{eq:matrix-product} give $\|B_m\|_\infty=1$ for all $m\leq N$. Also $e_0=0$ and $e_1\leq3\eta$.

Subtract the two matrix recurrences for $1\leq m<N$. Writing $\Delta_m=\widetilde B_m-B_m$ gives exactly
\begin{align*}
\widetilde a_m\Delta_{m+1}
={}&(\widetilde J-\widetilde b_m I)\Delta_m-\widetilde c_m\Delta_{m-1}\\
 &+\big(\widetilde J-J-(\widetilde b_m-b_m)I\big)B_m\\
 &-(\widetilde c_m-c_m)B_{m-1}-(\widetilde a_m-a_m)B_{m+1}.
\end{align*}
Since $\eta\leq1/(2N)$, $\widetilde a_m\geq1/(2N)$. The terms not containing a $\Delta$ have combined norm at most $6\eta$. Also $0\leq\widetilde b_m,\widetilde c_m\leq1$, so
\begin{equation}\label{eq:error-recursion}
e_{m+1}\leq4Ne_m+2Ne_{m-1}+12N\eta.
\end{equation}
Set $L=8N$. Induction gives $e_m\leq\eta L^m$ for $1\leq m\leq N$: the base case follows from $3\leq L$, and, for $m\geq1$,
\[
4N+\frac{2N}{L}+\frac{12N}{L^m}
\leq4N+\frac14+\frac32\leq8N=L.
\]
Thus every matrix entry differs from its reference by at most
\[
\eta(8N)^N=\frac12(s/4)^N.
\]
Lemma~\ref{lem:margin} gives the positive lower bound on triangle entries, while Lemma~\ref{lem:finite} makes all entries outside the triangle exactly zero. Finally $\widetilde R_m(1)=1$ follows by recurrence from the stochastic coefficient sums, so evaluating each product at $1$ gives coefficient sum one. Here $1$ belongs to the support because $\widetilde J\boldsymbol1=\boldsymbol1$ and Lemma~\ref{lem:finite} identifies that support with its spectrum. No product-positivity inequality for the perturbed system has been assumed.
\end{proof}

\begin{corollary}[Finite systems outside the geometric family]\label{cor:nonclassical}
For every $N\geq2$, take the Krawtchouk coefficients at $r=1/2$, and replace only
\[
a_1\longmapsto a_1-h,\qquad b_1\longmapsto b_1+h,
\qquad 0<h\leq\frac14(64N)^{-N}.
\]
These recurrences have nonnegative linearization and do not satisfy the constant-ratio hypothesis of Theorem~\ref{thm:criterion}. More generally, Theorem~\ref{thm:perturbation} admits a relatively open neighborhood of the asymmetric Krawtchouk array in the affine space defined by the row-sum and boundary equalities in \eqref{eq:basic}.
\end{corollary}
\begin{proof}
All initial and terminal conventions and stochastic sums are preserved. The bound on $h$ is half the radius \eqref{eq:radius} at $r=1/2$, and is smaller than $a_1=(N-1)/N$, so all required signs hold. The perturbation theorem applies. At index $1$, $b_1/c_1=1+2Nh>1$, whereas at index $N$ the ratio remains $1$; these indices differ because $N\geq2$. Thus this array lies outside the geometric family. At the asymmetric reference, all free coefficients constrained to be positive are strictly positive; the strict version of \eqref{eq:closeness} therefore contains a relatively open neighborhood in that affine space; the sign conditions define the admissible region within it.
\end{proof}

\section{Applications and scope}\label{sec:applications}
\begin{corollary}[Full positive Krawtchouk range]\label{cor:kraw}
For every $N\geq1$ and $0<p_0\leq1/2$, Krawtchouk polynomials orthogonal for binomial weights $\binom Nj p_0^{N-j}(1-p_0)^j$ have nonnegative linearization when normalized at $j=0$.
\end{corollary}
\begin{proof}
To fix the convention without relying on a named recurrence, put $r=p_0/(1-p_0)\in(0,1]$, and let $Y_1,\ldots,Y_N$ be independent variables with values $1,-r$ of probabilities $p_0,1-p_0$. Their mean is zero and $Y_i^2=(1-r)Y_i+r$. The normalized elementary symmetric sums
\[
K_i(j)=\binom Ni^{-1}\sum_{|S|=i}\prod_{l\in S}Y_l
\]
depend only on the number $j$ of negative coordinates. They are orthogonal, because the expectation of two distinct squarefree monomials is zero: some independent mean-zero variable occurs once. Their weights in $j$ are exactly the stated binomial weights. With $x=N^{-1}\sum_lY_l=1-j/(N(1-p_0))$, multiplication by $x$ gives
\[
xK_i=\frac{N-i}{N}K_{i+1}+(1-r)\frac iN K_i+r\frac iN K_{i-1}.
\]
Indeed, in multiplying the elementary sum by $\sum_lY_l$, indices outside a subset give $(i+1)$ copies of each subset of size $i+1$; the linear part of $Y_l^2$ gives $(1-r)i$ copies of each subset of size $i$; and the constant part gives $r(N-i+1)$ copies of each subset of size $i-1$. Division by the binomial normalizers gives the displayed recurrence, including $i=N$. Its leading coefficient in $x$ is positive, its degree is $i$, and its value at $j=0$ is one. These properties identify the normalized orthogonal polynomials for the binomial measure. Apply Theorem~\ref{thm:criterion} with $q=1$.
The range is sharp in this endpoint normalization: if $1/2<p_0<1$, then
$r=p_0/(1-p_0)>1$ and $g(1,1;1)=b_1=(1-r)/N<0$.
For $N=1$, this coefficient follows from the terminal relation.
\end{proof}

\begin{corollary}[The dual $q$-Krawtchouk case]\label{cor:additional}
For every $N\geq2$, $q>1$, and $0<r\leq1$, the recurrence \eqref{eq:parameters} has nonnegative linearization. In the normalization $R_0=1$, $R_1=x$, it is outside the Krawtchouk recurrence class.
\end{corollary}
\begin{proof}
Its increments are $d_i-d_{i-1}=q^{i-1}/[N]_q$, so all hypotheses hold. For a Krawtchouk recurrence in this normalization the increments of $b_i+c_i$ are constant. Here their ratio is $q>1$ and $N\geq2$, so they are not constant. An affine change of the support transporting the same endpoint-normalized basis leaves the multiplication coefficients in the basis unchanged; it cannot remove this distinction.
\end{proof}

\begin{proposition}[Relation to the existing algebraic construction]\label{prop:source}
The quadratic relations \eqref{eq:algebra} are those of Koelink~\cite[Theorem~3.1(ii), equation (3.3)]{koelink}, with $z_j=u(x^j)$ and the same parameters $p,q$. The recurrence in this paper is the normalized recurrence of his Theorem~7.1, with $r=1/p$; its polynomial-in-the-spectral-variable interpretation is the dual $q$-Krawtchouk interpretation described after his Theorem~7.8.
\end{proposition}
\begin{proof}
The source square relation is $u(x^j)^2=(q-1)u(x^j)\sum_{l>j}u(x^l)+(p-1)u(x^j)+pq^{N-j}$, which is \eqref{eq:algebra}. In Theorem~7.1 the coefficients of the next, current and previous indices are respectively $pq^i(1-q^{N-i})$, $(p-1)(1-q^i)$ and $1-q^i$. Division by $p(1-q^N)$ gives \eqref{eq:parameters} for $q>1$. For $0\leq f\leq N$, the corresponding spectral coordinate is
\[
x_f=\frac{p(1-q^{N-f})-(1-q^f)}{p(1-q^N)}
=\frac{q^{N-f}-rq^f+r-1}{q^N-1},\qquad r=1/p.
\]
Writing $\phi_f(w_d)$ for Koelink's spherical-function values, we have $R_d(x_f)=\phi_f(w_d)$ for $0\leq d\leq N$: the initial value is one and the normalized recurrences agree, so induction in $d$ identifies them.
As $q\to1$, dividing numerator and denominator by $q-1$ and using the finite geometric sums gives $x_f\to1-(1+r)f/N$.
The $q=1$ case follows directly from the construction above, without division by zero. His discussion after Theorem~7.8 explicitly interchanges polynomial degree and spectral index to obtain the dual family.
\end{proof}

The geometric criterion covers the full positive Krawtchouk range, including the symmetric endpoint. The perturbation criterion allows freely varying coefficients around each Krawtchouk array in the interior of the positive linearization range; its stated radius is conservative and decreases with $N$. Both hypotheses can be checked using a number of arithmetic comparisons proportional to $N$, without forming any general product coefficients or testing polynomial matrices. The proofs use positive algebra multiplication, a lower bound from a single positive term, and a norm estimate. They provide sufficient conditions, not a characterization of all finite systems with property (P). The geometric parameter family and underlying algebra predate Szwarc's survey; this paper makes no claim of historical priority for these constructions or criteria.

\section{Szwarc's quadratic system}\label{sec:szwarc}

Problem~2 of~\cite[\S5]{szwarc} does not stop at the Krawtchouk recurrence. With his parameter written as $p_0$, Szwarc introduces the symmetric recurrence of length $2N+1$
\begin{equation}\label{eq:szwarc-q}
yQ_{2n}=(N-n)Q_{2n+1}+nQ_{2n-1},\qquad
yQ_{2n-1}=(1-p_0)Q_{2n}+p_0Q_{2n-2},
\end{equation}
records that $Q_{2n}(y)$ is the Krawtchouk polynomial $K_n(y^2)$ of his recurrence, states that the $Q_n$ have nonnegative linearization for $p_0\leq1/2$, and asks for a general criterion that would handle this case. His Theorem~3.3 for infinite symmetric systems does not apply. In the normalization \eqref{eq:szwarc-q} its monotonicity hypotheses hold, and its hypothesis $\alpha_n\leq\gamma_n$ holds at odd indices exactly when $p_0\leq1/2$; at an even index $2n$ it reads $n\leq N-n$ and fails in the upper half. Some failure of this kind is forced in every finite system, whose terminal upward coefficient vanishes. This section proves the positivity of \eqref{eq:szwarc-q} by a graded version of the algebra of Section~\ref{sec:algebra}, and then gives an explicit coefficient criterion for symmetric finite systems near \eqref{eq:szwarc-q}.

Put $M=2N$, $R_{2n}(t)=Q_{2n}(\sqrt Nt)$ and $R_{2n-1}(t)=\sqrt N\,Q_{2n-1}(\sqrt Nt)$. Substitution in \eqref{eq:szwarc-q} shows that the $R_k$ satisfy \eqref{eq:recurrence}, with $N$ replaced by $M$, for the array
\begin{equation}\label{eq:qarray}
b_k=0,\qquad
(a_{2n},c_{2n})=\Big(\frac{N-n}{N},\frac nN\Big),\qquad
(a_{2n-1},c_{2n-1})=(1-p_0,\,p_0).
\end{equation}
It satisfies \eqref{eq:basic} with $M$ in place of $N$ whenever $0<p_0<1$. The two normalizations differ by positive factors and a positive dilation, so they have the same product coefficients up to positive factors.

\begin{theorem}[Szwarc's quadratic system]\label{thm:szwarc}
Let $N\geq1$ and $0<p_0\leq1/2$. The array \eqref{eq:qarray} has a symmetric positive orthogonality measure on $2N+1$ points, and for all $0\leq m,n\leq 2N$
\[
R_mR_n=\sum_{k=0}^{2N}g(m,n;k)R_k,\qquad g(m,n;k)\geq0,\qquad\sum_k g(m,n;k)=1
\]
on that support. If $1/2<p_0<1$, then $g(2,2;2)<0$.
\end{theorem}

Throughout this section put $r=p_0/(1-p_0)$, so that $0<r\leq1$ exactly when $p_0\leq1/2$. Let $A$ be the commutative algebra generated by $y_1,\ldots,y_N$ subject to $y_j^2=(1-r)y_j+r$. This is the algebra of Section~\ref{sec:algebra} with $q=1$ after the rescaling $y_j=z_j/p$ used in Lemma~\ref{lem:margin}. Its squarefree monomials $y_S$ form a basis, and for $r\leq1$ all multiplication coefficients are nonnegative. Call an element of $A$ nonnegative if its coefficients in this basis are nonnegative.

\begin{lemma}[Edge idempotents]\label{lem:edge}
For $1\leq l\leq N$ put $f_l=(r+y_l)/(1+r)$. Then $y_lf_l=f_l$ and $f_l^2=f_l$. The ideal $f_lA$ has basis $f_ly_S$, $S\subseteq\{1,\ldots,N\}\setminus\{l\}$, and $f_l u$ has nonnegative coefficients in this basis whenever $u\in A$ is nonnegative.
\end{lemma}
\begin{proof}
First $(1+r)y_lf_l=ry_l+(1-r)y_l+r=(1+r)f_l$. Then $(1+r)f_l^2=rf_l+y_lf_l=(1+r)f_l$. Let $A'$ be generated by the $y_j$ with $j\neq l$, so that $A=A'\oplus y_lA'$. Since $f_ly_l=f_l$, $f_lA=f_lA'$. The map $u\mapsto f_lu=(ru+y_lu)/(1+r)$ is injective on $A'$ by the direct sum, which gives the basis. For the last claim, $f_ly_S=f_ly_{S\setminus\{l\}}$ for every $S$.
\end{proof}

Let $V=A\oplus\bigoplus_{l=1}^Nf_lA$, on which $A$ acts by multiplication in every summand, and call a vector nonnegative if all of its coordinates are nonnegative in the bases above. Define the $A$-linear operator
\[
Y(u;v_1,\ldots,v_N)=\Big(\frac1N\sum_{l=1}^Nv_l\,;\,f_1u,\ldots,f_Nu\Big).
\]
Put $U_{2n}=\binom Nn^{-1}\sum_{|S|=n}y_S\in A$ for $0\leq n\leq N$. For $0\leq n\leq N-1$, let $U_{2n+1}$ be the vector in $\bigoplus_lf_lA$ whose $l$th coordinate is $\binom{N-1}n^{-1}\sum f_ly_S$, summed over $S\subseteq\{1,\ldots,N\}\setminus\{l\}$ with $|S|=n$.

\begin{lemma}[Three-term action]\label{lem:qthree}
With $U_{-1}=U_{2N+1}=0$ and the coefficients \eqref{eq:qarray}, $YU_k=a_kU_{k+1}+c_kU_{k-1}$ for $0\leq k\leq2N$.
\end{lemma}
\begin{proof}
The $l$th coordinate of $Y\sum_{|S|=n}y_S$ is $\sum_{S\not\ni l}f_ly_S+\sum_{S'\not\ni l}f_ly_{S'}$, with $|S|=n$ and $|S'|=n-1$, because $f_ly_S=f_ly_{S\setminus\{l\}}$. Dividing by $\binom Nn$ and using $\binom{N-1}n/\binom Nn=(N-n)/N$ and $\binom{N-1}{n-1}/\binom Nn=n/N$ gives the even case. For the odd case, $(1+r)f_ly_S=ry_S+y_{S\cup\{l\}}$ when $l\notin S$. Each $n$-set occurs for $N-n$ choices of $l$ and each $(n+1)$-set for $n+1$ choices, so
\[
YU_{2n+1}=\frac{r(N-n)\binom Nn U_{2n}+(n+1)\binom N{n+1}U_{2n+2}}{N(1+r)\binom{N-1}n}
=\frac{rU_{2n}+U_{2n+2}}{1+r}.
\]
Finally $r/(1+r)=p_0$ and $1/(1+r)=1-p_0$.
\end{proof}

\begin{proof}[Proof of Theorem~\ref{thm:szwarc}]
Lemma~\ref{lem:finite} gives a positive measure on the $2N+1$ zeros of the terminal polynomial $T(t)=tR_{2N}(t)-R_{2N-1}(t)$. Since $b\equiv0$, the matrix $J$ is conjugate to $-J$ by $\operatorname{diag}((-1)^i)$, which fixes $e_0$; hence the measure is symmetric and $R_k(-t)=(-1)^kR_k(t)$. Put $\Phi(f)=f(Y)U_0$, where $U_0=(1;0,\ldots,0)$. Lemma~\ref{lem:qthree} gives $\Phi(R_k)=U_k$ for $k\leq2N$ and $\Phi(T)=YU_{2N}-U_{2N-1}=0$. Every polynomial vanishing on the support is a multiple of $T$, so $\Phi$ maps it to zero. Hence the product identity on the support becomes
\[
R_m(Y)R_n(Y)U_0=\sum_kg(m,n;k)U_k.
\]
The supports of the $U_k$ in the bases of $V$ are nonempty and disjoint: $U_{2n}$ uses the $y_S$ with $|S|=n$, and $U_{2n+1}$ uses the $f_ly_S$ with $l\notin S$, $|S|=n$. Each basis vector lies in exactly one of them, with positive coefficient. Consequently $g(m,n;k)\geq0$ for every $k$ whenever the left side is nonnegative.

Suppose $m$ is even. Since polynomials in $Y$ commute and $Y$ is $A$-linear,
\[
R_m(Y)R_n(Y)U_0=R_n(Y)(U_mU_0)=U_m\,R_n(Y)U_0=U_mU_n.
\]
Here $U_m\in A$ and $U_n$ are nonnegative, so their product is nonnegative by the positivity of $A$ and Lemma~\ref{lem:edge}. The same holds when $n$ is even, since $g(m,n;k)=g(n,m;k)$. If $m,n$ are odd and $k$ is odd, then $g(m,n;k)=0$ by symmetry of the measure. If $m,n$ are odd and $k$ is even, then $g(m,n;k)\int R_k^2\,d\mu=\int R_mR_nR_k\,d\mu=g(k,m;n)\int R_n^2\,d\mu\geq0$. The point $t=1$ belongs to the support because $J\boldsymbol1=\boldsymbol1$, and $R_k(1)=1$, so the coefficients sum to one.

For $1/2<p_0<1$ we have $r>1$. The identities above are purely algebraic and remain valid, but $A$ is no longer positive. In particular $U_2^2=N^{-2}\big((1-r)\sum_jy_j+Nr+2\sum_{i<j}y_iy_j\big)$, whose $U_2$ coefficient $g(2,2;2)=(1-r)/N$ is negative.
\end{proof}

Since $Y^2$ acts on $A$ by multiplication by $N^{-1}\sum_lf_l=(r+X)/(1+r)$, with $X=N^{-1}\sum_ly_l$, the even part of the theorem is Corollary~\ref{cor:kraw}: $R_{2n}(t)=K_n\big((1+r)t^2-r\big)$ in the normalization of that corollary. Thus Theorem~\ref{thm:szwarc} recovers the full positive Krawtchouk range by Szwarc's quadratic route. At $p_0=1/2$ the operator $Y$ is the normalized adjacency operator of the subdivided $N$-cube seen from a vertex, and the theorem reduces to counting in that graph.

\begin{lemma}[Margin for the quadratic system]\label{lem:qmargin}
Let $0<p_0<1/2$ and $s=\min(r,1-r)$. For the array \eqref{eq:qarray},
\[
g(m,n;k)\geq(rs/8)^N\qquad\text{whenever }|m-n|\leq k\leq m+n\text{ and }m+n+k\text{ is even}.
\]
All other coefficients vanish.
\end{lemma}
\begin{proof}
The vanishing statements are Lemma~\ref{lem:finite} and the symmetry used above. We use $\binom Nk\leq2^N$ throughout. Up to the order of $m,n$, there are three cases.

If $m,n,k$ are even, the coefficient is a Krawtchouk coefficient computed in $A$, so Lemma~\ref{lem:margin} gives the lower bound $(s/4)^N\geq(rs/8)^N$.

Let $m=2a$, $n=2b+1$ and $k=2c+1$, so that
\[
\max(a-b-1,\,b-a)\leq c\leq\min(a+b,\,N-1).
\]
Fix $l$ and a set $U\not\ni l$ with $|U|=c$. The $l$th coordinate of $U_mU_n$ is
\[
\binom Na^{-1}\binom{N-1}b^{-1}\sum_{|T|=a}\ \sum_{S\not\ni l,\ |S|=b}f_l\,y_{T\setminus\{l\}}\,y_S .
\] Every monomial coefficient of $y_{T\setminus\{l\}}y_S$ is a product of at most $N-1$ factors from $\{r,1-r\}$, and is therefore at least $s^{N-1}$. Put $a'=|T\setminus\{l\}|$, which can equal $a-1$ when $a\geq1$ and $a$ when $a\leq N-1$. In the $N-1$ coordinates other than $l$, sets of sizes $a'$ and $b$ with overlap $o$, retaining the linear term on $t\leq o$ overlap coordinates, give monomials of size $a'+b-2o+t$. As $o$ ranges over its admissible values these sizes fill the interval $\big[|a'-b|,\min(a'+b,N-1)\big]$. The two admissible values of $a'$ cover every $c$ above. Hence $f_ly_U$ has coefficient at least $s^{N-1}\binom Na^{-1}\binom{N-1}b^{-1}$, and comparison with $U_k$ gives
\[
g(m,n;k)\geq\binom{N-1}c\frac{s^{N-1}}{\binom Na\binom{N-1}b}\geq\frac{2s^{N-1}}{4^N}\geq\Big(\frac s4\Big)^N.
\]

Finally let $m=2b+1$, $n=2c+1$ and $k=2a$. With $\pi_i$ as in the proof of Lemma~\ref{lem:finite}, \eqref{eq:qarray} gives $\pi_{2a}=\binom Nar^{-a}$ and $\pi_{2c+1}=\binom{N-1}cr^{-c}/p_0$. Orthogonality gives
\[
g(2b+1,2c+1;2a)=g(2a,2b+1;2c+1)\frac{\pi_{2a}}{\pi_{2c+1}}
\geq\frac{2s^{N-1}}{4^N}\cdot\frac{p_0r^N}{2^{N-1}}
\geq\frac{2s^{N-1}r^{N+1}}{8^N},
\]
using $r^{c-a}\geq r^N$ and $p_0\geq r/2$. Since $s\leq r$, the last quantity is at least $(rs/8)^N$.
\end{proof}

The proof of Theorem~\ref{thm:perturbation} uses only three properties of its reference: nonnegative linearization with sums one, a lower bound for the upward coefficients, and a margin on the known support. We record it in that generality.

\begin{lemma}[Stability of a reference with a margin]\label{lem:stability}
Let $a,b,c$ satisfy \eqref{eq:basic} on the indices $0,\ldots,M$, have nonnegative linearization, and satisfy $a_i\geq\alpha$ for $i<M$, where $0<\alpha\leq1$. Suppose $g(m,n;k)\geq\delta>0$ on a set $\Sigma$ of index triples. Let $\widetilde a,\widetilde b,\widetilde c$ satisfy \eqref{eq:basic}, with $\widetilde b_i\geq0$, and differ from $a,b,c$ by at most $\eta$ in every entry. If $\eta\leq\alpha/2$ and $\eta(8/\alpha)^M\leq\delta/2$, then $|\widetilde g(m,n;k)-g(m,n;k)|\leq\delta/2$ for all $m,n,k$; in particular $\widetilde g\geq\delta/2$ on $\Sigma$.
\end{lemma}
\begin{proof}
Repeat the proof of Theorem~\ref{thm:perturbation} with $N$ replaced by $M$. Now $\widetilde a_m\geq\alpha/2$, so the subtracted recurrences give $e_{m+1}\leq(4e_m+2e_{m-1}+12\eta)/\alpha$. With $L=8/\alpha\geq8$, induction gives $e_m\leq\eta L^m$: the base case is $e_1\leq3\eta$, and the step needs $4L^m+2L^{m-1}+12\leq8L^m$, which holds for $m\geq1$. Every entry of $\widetilde B_m-B_m$ is therefore at most $\eta L^M\leq\delta/2$, and \eqref{eq:matrix-product} identifies these entries with the product coefficients.
\end{proof}

Theorem~\ref{thm:perturbation} is the case $M=N$, $\alpha=1/N$ and $\delta=(s/4)^N$ of this lemma.

\begin{theorem}[Symmetric criterion near Szwarc's system]\label{thm:symmetric}
Let $N\geq1$, $0<p_0<1/2$, $r=p_0/(1-p_0)$, $s=\min(r,1-r)$, and
\[
\eta=\frac12\Big(\frac{rs}8\Big)^N(16N)^{-2N}.
\]
Let $c_k$ be the downward coefficients in \eqref{eq:qarray}. Take real numbers $\widetilde c_1,\ldots,\widetilde c_{2N-1}$ with $|\widetilde c_k-c_k|\leq\eta$, and put $\widetilde c_0=0$, $\widetilde c_{2N}=1$, $\widetilde b_k=0$ and $\widetilde a_k=1-\widetilde c_k$. The resulting symmetric finite system has nonnegative linearization with sums one. Moreover $\widetilde g(m,n;k)\geq\frac12(rs/8)^N$ whenever $|m-n|\leq k\leq m+n$ and $m+n+k$ is even, and $\widetilde g(m,n;k)=0$ otherwise.
\end{theorem}
\begin{proof}
The perturbed array satisfies \eqref{eq:basic} with $M=2N$: $\widetilde a_k\geq1/N-\eta>0$ for even $k<2N$, $\widetilde a_k\geq1/2-\eta>0$ for odd $k$, and $\widetilde c_k\geq\min(p_0,1/N)-\eta>0$ for $k>0$. Theorem~\ref{thm:szwarc} shows that the reference has nonnegative linearization, and $a_k\geq1/(2N)$ for $k<2N$. Lemma~\ref{lem:qmargin} gives the margin $\delta=(rs/8)^N$ on the stated set. Apply Lemma~\ref{lem:stability} with $\alpha=1/(2N)$: then $\eta\leq\alpha/2$ and $\eta(16N)^{2N}=\delta/2$. Off the stated set, $\widetilde g$ vanishes by Lemma~\ref{lem:finite} and by the symmetry of the perturbed measure, which holds because $\widetilde b\equiv0$. The sums equal one because $\widetilde R_m(1)=1$ and $1$ lies in the support.
\end{proof}

\begin{corollary}[Symmetric systems outside both families]\label{cor:symmetric-new}
For $N\geq2$, $0<p_0<1/2$ and $0<h\leq\eta$, the array obtained from \eqref{eq:qarray} by $c_1\mapsto p_0+h$, $a_1\mapsto1-p_0-h$ has nonnegative linearization. It is neither of the form \eqref{eq:qarray} nor of the form \eqref{eq:parameters}. More generally, Theorem~\ref{thm:symmetric} admits a relatively open set of symmetric stochastic arrays, with $2N-1$ free coefficients, around each array \eqref{eq:qarray} with $0<p_0<1/2$.
\end{corollary}
\begin{proof}
Theorem~\ref{thm:symmetric} applies. The odd downward coefficients $c_1=p_0+h$ and $c_3=p_0$ differ, so the array is not of the form \eqref{eq:qarray}; here $N\geq2$ ensures $3<2N$. In \eqref{eq:parameters}, on $2N+1$ indices, the downward coefficients $c_i=r[i]_q/[2N]_q$ are strictly increasing, whereas $c_3<c_1$.
\end{proof}

The quadratic map of~\cite[\S3.5]{szwarc} transfers these conclusions to nonsymmetric systems. If a symmetric finite system has nonnegative linearization, then its even polynomials, written in the variable $t^2$, form a finite system with nonnegative linearization: products of even polynomials involve only even polynomials. Theorem~\ref{thm:symmetric} therefore also yields nonsymmetric finite systems with nonnegative linearization, namely the quadratic transforms of the arrays it admits; their recurrence coefficients are given by the formula of~\cite[\S3.5]{szwarc}.

\section{Open problems}\label{sec:open}

\paragraph{Inequalities beyond geometric increments.}
The geometric criterion imposes exact equalities, whereas the perturbation criterion provides small neighborhoods. Find direct inequalities on successive recurrence coefficients that cover larger nonclassical regions while retaining the full Krawtchouk range. The terminating row and products of degree sum greater than $N$ must remain part of the conclusion.

\paragraph{Sharper stability radii.}
For fixed $N$ and $0<r<1$, determine or estimate the largest coefficient radius for which every admissible stochastic perturbation has property (P). The bound in Theorem~\ref{thm:perturbation} comes from a single-term lower margin and a worst-case recurrence amplification. Both estimates discard substantial information. A useful improvement would retain a transparent coefficient test and describe its dependence on both $N$ and the bias parameter.

\paragraph{The symmetric endpoint.}
At $r=1$, the reference margin in Lemma~\ref{lem:margin} vanishes because parity forces many coefficients to be zero. Symmetric perturbations preserve the parity zeros, and Theorem~\ref{thm:symmetric} exploits this for Szwarc's quadratic system when $p_0<1/2$. At $p_0=1/2$, and for the symmetric Krawtchouk array, further coefficients on the parity support vanish. Determine perturbation directions and structural constraints that preserve positivity there; continuity alone does not control coefficients emerging from those zeros.

\paragraph{A monotonicity criterion for finite systems.}
Szwarc's Theorems~3.2 and~3.3 in~\cite{szwarc} are open conditions of monotonicity type for infinite systems. For \eqref{eq:szwarc-q} all their hypotheses except $\alpha_n\leq\gamma_n$ at the upper even indices hold. Find inequality hypotheses of this kind, valid for finite systems including the terminal indices, that cover Theorem~\ref{thm:szwarc} without the exact algebraic structure used here.

\paragraph{Other finite families.}
Find analogous coefficient regions around other finite orthogonal systems and compare their geometry with the Krawtchouk neighborhoods. An algebra with a positive basis may supply the reference margin, but a recurrence criterion must express the resulting admissible deformations without requiring the unknown product coefficients as input.

\section*{AI-use disclosure}

OpenAI Codex was used to develop the proof arguments, search and check source passages, construct exact-arithmetic diagnostics, draft and revise this manuscript, and prepare the Lean translation. Separate AI-assisted reviews checked the algebraic proof, boundary cases, normalizations, and quantitative estimates; these reviews are not human peer review. Anthropic Claude was used to analyse Szwarc's Problem~2 against the survey text, to explore discrete hyperbolic certificates numerically, and to develop and draft Section~\ref{sec:szwarc} together with its exact diagnostics; that section has not been independently reviewed. The author set the research objective and retains responsibility for publication decisions and the final mathematical claims. No independent human verification is asserted here. Statements about machine verification refer only to the specific compiled Lean declarations and their documented dependencies, not to the presence of a Lean project or to numerical agreement.

\section*{Code and formalization}

The companion source tree contains \path{probe.py} and \path{verify.py}. The command
\begin{quote}\verb|python3 -u verify.py --max-n 7|\end{quote}
compares recurrence-matrix multiplication with direct multiplication modulo the terminal polynomial on exact rational inputs, checks the perturbation lower bounds, and replays two rejected coefficient-profile tests. It also checks Theorem~\ref{thm:szwarc}, the margin of Lemma~\ref{lem:qmargin} and the negative coefficient for $p_0>1/2$ on exact Szwarc arrays, the identities of Lemma~\ref{lem:qthree} in the graded model, and the lower bound of Theorem~\ref{thm:symmetric} on rational symmetric perturbations. These computations are diagnostics and do not replace the proofs at arbitrary size.

The Lean development formalizes the two main criteria as follows:
\begin{center}
\begin{tabular}{@{}p{0.36\linewidth}p{0.59\linewidth}@{}}
Paper result & Lean declaration\\[3pt]
Theorem~\ref{thm:criterion}: geometric criterion & \path{Szkraw.geometric_criterion}\\[3pt]
Theorem~\ref{thm:perturbation}: perturbation & \path{Szkraw.perturbation_criterion}
\end{tabular}
\end{center}
The first includes a positive finite orthogonality measure, distinct support points, the terminal relation, and all finite product identities. The second includes the explicit radius, positive triangle margin, and exact zeros outside the triangle. Both have compiled proofs without added axioms or admitted goals. The results of Section~\ref{sec:szwarc} have not been formalized.

The companion code is available at \url{https://github.com/skypher/szkraw}; the fixed revision cited in~\cite{code} identifies the reviewed proof and audit code. Run \verb|make lean| to build the development and \verb|make axioms| to build and audit the principal statements. The latter invokes \path{lean/CheckAxioms.lean}, which rejects dependencies outside \path{propext}, \path{Classical.choice}, and \path{Quot.sound} and tests rejection of unexpected axioms. The toolchain and dependencies are pinned in \path{lean/lean-toolchain} and \path{lean/lake-manifest.json}. The declaration-by-declaration correspondence, auxiliary formalization coverage, and verification commands are recorded in \path{RESEARCH.md} at that revision. Historical attribution and the open problems are prose claims, not machine-verified statements.

\begin{thebibliography}{9}
\bibitem{askey} R. Askey, \emph{Three notes on orthogonal polynomials}, Stichting Mathematisch Centrum, Afdeling Toegepaste Wiskunde, report TW~114/69, Amsterdam, October 1969; Part~II: \emph{Linearization of the product of orthogonal polynomials II}, \url{https://ir.cwi.nl/pub/7741}.
\bibitem{eagleson} G. K. Eagleson, \emph{A characterization theorem for positive definite sequences on the Krawtchouk polynomials}, Australian Journal of Statistics \textbf{11} (1969), no.~1, 29--38, \href{https://doi.org/10.1111/j.1467-842X.1969.tb00004.x}{\nolinkurl{doi:10.1111/j.1467-842X.1969.tb00004.x}}.
\bibitem{koelink} H. T. Koelink, \emph{$q$-Krawtchouk polynomials as spherical functions on the Hecke algebra of type B}, Transactions of the American Mathematical Society \textbf{352} (2000), no.~10, 4789--4813, \href{https://doi.org/10.1090/S0002-9947-00-02588-5}{\nolinkurl{doi:10.1090/S0002-9947-00-02588-5}}; preprint arXiv:q-alg/9605040v1 (1996), \url{https://arxiv.org/abs/q-alg/9605040v1}. Theorem and equation locators refer to the preprint.
\bibitem{miclo} L. Miclo, \emph{On the hypergroup property}, Annales de la Facult\'e des Sciences de Toulouse, s\'erie~6, \textbf{26} (2017), no.~2, 417--435, \href{https://doi.org/10.5802/afst.1539}{\nolinkurl{doi:10.5802/afst.1539}}.
\bibitem{szwarc} R. Szwarc, \emph{Orthogonal polynomials and Banach algebras}, in \emph{Inzell Lectures on Orthogonal Polynomials}, W. zu Castell, F. Filbir, and B. Forster (eds.), Advances in the Theory of Special Functions and Orthogonal Polynomials, vol.~2, Nova Science Publishers, 2005, 103--139; \S5, Problem~2. Author version (2004, 97--133; Problem~2 on printed p.~125): \url{https://www.math.uni.wroc.pl/~szwarc/pdf/inzell.pdf}.
\bibitem{szwarc2003} R. Szwarc, \emph{A necessary and sufficient condition for nonnegative product linearization of orthogonal polynomials}, Constructive Approximation \textbf{19} (2003), 565--573, \href{https://doi.org/10.1007/s00365-002-0524-z}{\nolinkurl{doi:10.1007/s00365-002-0524-z}}.
\bibitem{code} L. P. Polzer, \emph{Companion code for Coefficient criteria for finite nonnegative linearization} (2026), revision \texttt{e65b421} (full identifier in the URL), \url{https://github.com/skypher/szkraw/tree/e65b421f25cec90c8d4097363da80386592ce5ed}.
\end{thebibliography}
\end{document}
